\documentclass[11pt]{article}
\usepackage[a4paper,margin=1in]{geometry}
\usepackage{fontspec}
\usepackage[boldfont=false]{xeCJK}
\setCJKmainfont{FandolSong-Regular.otf}
\usepackage{amsmath}
\usepackage{amssymb}
\usepackage{array}
\usepackage{booktabs}
\usepackage{graphicx}
\usepackage{adjustbox}
\usepackage{enumitem}
\setlist[itemize]{topsep=2pt,partopsep=0pt,parsep=0pt,itemsep=2pt,leftmargin=1.4em}
\setlist[enumerate]{topsep=2pt,partopsep=0pt,parsep=0pt,itemsep=2pt,leftmargin=1.6em}
\usepackage{parskip}
\usepackage{fvextra}
\fvset{breaklines=true,breakanywhere=true,fontsize=\small}
\setlength{\parindent}{0pt}

\begin{document}

\begin{center}
  {\LARGE\bfseries Respiration}\\[0.35em]
  {\large Igcse Biology}
\end{center}
\vspace{0.6em}

\section*{Respiration and energy}

\textbf{Respiration} 呼吸作用 is the release of \textbf{energy} 能量 from food, and it happens in every living cell, all the time. (It is \textbf{not} the same as breathing.) Cells use this energy for many jobs:

\begin{itemize}
\item \textbf{muscle} 肌肉 contraction (movement).

\item \textbf{protein synthesis} 蛋白质合成 (building proteins).

\item cell \textbf{division} 分裂 (making new cells).

\item \textbf{active transport} 主动运输.

\item growth.

\item the passage of \textbf{nerve impulses} 神经冲动.

\item keeping a constant body \textbf{temperature} 温度.

\end{itemize}

You can investigate respiration in \textbf{yeast} 酵母: warmer yeast respires faster (up to its best temperature), giving off bubbles of \textbf{carbon dioxide} 二氧化碳.

\section*{Aerobic respiration}

\textbf{Aerobic respiration} 有氧呼吸 uses \textbf{oxygen} 氧气 to break down \textbf{nutrient molecules} 营养物质 (mainly \textbf{glucose} 葡萄糖) and release energy.

Word equation:

\[
\text{glucose} + \text{oxygen} \rightarrow \text{carbon dioxide} + \text{water}
\]

\textbf{(Supplement)} Balanced chemical equation:

\[
\text{C}_6\text{H}_{12}\text{O}_6 + 6\text{O}_2 \rightarrow 6\text{CO}_2 + 6\text{H}_2\text{O}
\]

Aerobic respiration releases a \textbf{lot} of energy from each glucose molecule.

\section*{Anaerobic respiration}

\textbf{Anaerobic respiration} 无氧呼吸 breaks down glucose to release energy \textbf{without} oxygen. It releases \textbf{much less} energy from each glucose molecule than aerobic respiration, because the glucose is not fully broken down.

\subsection*{In yeast}

\[
\text{glucose} \rightarrow \text{alcohol} + \text{carbon dioxide}
\]

The \textbf{alcohol} 酒精 made by yeast is used to make bread rise and to brew drinks.

\subsection*{In muscles (Supplement)}

During hard exercise your muscles cannot get enough oxygen, so they respire anaerobically:

\[
\text{glucose} \rightarrow \text{lactic acid}
\]

The \textbf{lactic acid} 乳酸 builds up in the muscles and the blood. This creates an \textbf{oxygen debt} 氧债 — the extra oxygen the body will need later to break that lactic acid down.

After you stop exercising, the oxygen debt is repaid:

\begin{itemize}
\item your \textbf{heart rate} 心率 stays high, carrying the lactic acid in the blood from the muscles to the \textbf{liver} 肝脏.

\item you keep breathing deeply and quickly, taking in extra oxygen.

\item the liver uses this oxygen to break the lactic acid down by aerobic respiration.

\end{itemize}

\section*{Exam tips}

\begin{itemize}
\item Respiration releases energy in \textbf{all} living cells, all the time — it is not breathing.

\item Aerobic: glucose + oxygen → carbon dioxide + water, and it gives \textbf{a lot} of energy. Learn the balanced equation too.

\item Anaerobic (no oxygen) gives \textbf{much less} energy. In yeast: glucose → alcohol + carbon dioxide. In muscles: glucose → lactic acid.

\item Lactic acid causes an \textbf{oxygen debt}, repaid by a fast heart rate and deep breathing, with the liver breaking the lactic acid down.

\end{itemize}


\end{document}
